% Figure: the power-angle curve of a two-bus AC link. Real power transfer
% P = (V1 V2 / X) sin(delta) rises with the voltage-angle difference delta
% to a maximum at 90 degrees and falls beyond, the static stability limit.
% A horizontal operating-power line cuts the curve at a stable point on the
% rising branch and an unstable point on the falling branch. Values
% illustrative (per-unit), with V1 = V2 = 1.0 pu and reactance X = 0.2 pu so
% the peak transfer is 1/0.2 = 5 pu.
\begin{figure}[htbp]
\centering
\begin{tikzpicture}
\begin{axis}[
  width=12.5cm, height=7.5cm,
  xlabel={voltage-angle difference $\delta$ (degrees)},
  ylabel={real power transfer $P$ (pu)},
  xmin=0, xmax=180, ymin=0, ymax=5.6,
  xtick={0,30,60,90,120,150,180},
  ytick={0,1,2,3,4,5},
  axis lines=left,
  tick label style={font=\scriptsize},
  label style={font=\small},
  legend style={font=\scriptsize, at={(0.97,0.03)}, anchor=south east, draw=none},
]
% power-angle curve P = 5 sin(delta)
\addplot[engblue, very thick, domain=0:180, samples=181] {5*sin(x)};
\addlegendentry{$P=(V_1V_2/X)\sin\delta$}
% operating power line at P = 3.5 pu
\addplot[engred, thick, dashed, domain=0:180, samples=2] {3.5};
\addlegendentry{operating power $=3.5$ pu}
% stable operating point: 5 sin(delta) = 3.5 -> delta = 44.4 deg
\addplot[engorange, only marks, mark=*, mark size=2.2pt] coordinates {(44.4,3.5)};
\node[anchor=west, font=\scriptsize, color=engorange] at (axis cs:46,3.0) {stable};
% unstable operating point: delta = 135.6 deg
\addplot[enggray, only marks, mark=o, mark size=2.2pt] coordinates {(135.6,3.5)};
\node[anchor=east, font=\scriptsize, color=enggray] at (axis cs:133,3.0) {unstable};
% peak / static stability limit at 90 deg
\addplot[englime, only marks, mark=triangle*, mark size=2.5pt] coordinates {(90,5)};
\node[anchor=south, font=\scriptsize, color=englime] at (axis cs:90,5.05) {limit $P_{\max}=V_1V_2/X$};
\end{axis}
\end{tikzpicture}
\caption[The power-angle curve of a two-bus AC link]{The power-angle curve
of a lossless two-bus AC link: real power transfer follows
$P=(V_1V_2/X)\sin\delta$ in the voltage-angle difference $\delta$ across the
connecting reactance $X$. Transfer rises to a maximum at $\delta=90^\circ$,
the static stability limit $P_{\max}=V_1V_2/X$, and falls beyond it. A given
operating power (red) is met at two angles: a stable point on the rising
branch, where a small increase in load draws a restoring increase in
transfer, and an unstable point on the falling branch, where the same
increase reduces transfer and the machines lose synchronism. Operating
margin is the gap between the working angle and $90^\circ$. Per-unit values
illustrative, with $V_1=V_2=1.0$ pu and $X=0.2$ pu.}
\label{fig:vol04ch14:power-angle}
\end{figure}
